% !PRINTED 389
\weberpage{389}
%<<CONT
\gesperrt{Zahlen ausschließen, die im Zähler oder im Nenner nicht relativ prim zu $S$ sind.} (Im Falle der Ordnungen müssen die Primfaktoren des Führers $Q$ in $S$ enthalten sein.)

Wir bilden nun aus $M$ einen Teiler $A$, in den wir alle Zahlen $a$ aufnehmen, die einer Kongruenz

\weq{2}{N(\omega) \equiv a \pmod{m}}

genügen, d.\,h. für die eine Zahl $\omega$ in $O$ existiert, deren Norm mit $a$ kongruent ist.

1. \gesperrt{Diese Zahlen heißen Normenreste für den Modul $m$.}

Aus der Definition der Normenreste ergibt sich, daß die Normenreste der Multiplikation und Division gegenüber eine Gruppe bilden.

Ferner ergibt sich, \gesperrt{daß jedes Quadrat und jeder quadratische Rest (modulo $m$) zugleich Normenrest ist, daß also jeder Normennichtrest zugleich quadratischer Nichtrest ist.} Aus der Gruppennatur von $A$ folgt:

\gesperrt{Das Produkt aus Normenrest und Normenrest ist Normenrest.}

\gesperrt{Das Produkt aus Normenrest und Normennichtrest ist Normennichtrest.}

Die Gruppe der Normenreste hat stets einen endlichen Index $(Z, A)$, denn sie vereinigt in sich ganze \gesperrt{Zahlklassen} nach dem Modul $m$, deren Anzahl endlich ist.

\S\,107. \gesperrt{Zusammensetzung der Normenrestgruppen.}

Wir wollen zwei Moduln $m_1, m_2$ betrachten, die zueinander relativ prim sind, und bezeichnen die Gruppen der Normenreste von $m_1$ und $m_2$ mit $A_1$ und $A_2$. Setzen wir für den Augenblick

\[ (Z, A_1) = \mu \]

und zerlegen

\weq{1}{Z = a_1 A_1 + a_2 A_1 + \cdots + a_\mu A_1,}

so können wir die $a_1, a_2, \ldots, a_\mu$ durch beliebige nach dem Modul $m_1$ kongruente Zahlen ersetzen. Da $m_1$ und $m_2$ relativ prim vorausgesetzt sind, so lassen sich, welches auch die rationalen Zahlen $c_i$ seien, rationale $x_i$ aus den Kongruenzen

\[ a_i + x_i m_1 \equiv c_i \pmod{m_2} \]

bestimmen, und man kann also die $a_i$ so wählen, daß sie in $A_2$ enthalten sind. Man kann also jede Zahl $z$ in $Z$ als Produkt
%CONT>>
